\chapter*{منابع و مآخذ}
\addcontentsline{toc}{chapter}{منابع و مآخذ}
\markboth{منابع و مآخذ}{منابع و مآخذ}

\begin{flushright}
\textit{کتاب‌شناسی منتخب برای مطالعهٔ بیشتر}
\end{flushright}

\vspace{10pt}

\section*{منابع فارسی}

\begin{itemize}[label=--]
    \item رابرت هایلبرونر، \textit{بزرگان اقتصاد}، ترجمه احمد شهسا، انتشارات علمی و فرهنگی.
    \item کوئنتین اسکینر، \textit{بنیادهای اندیشه سیاسی مدرن}، ترجمه عباس کاظم‌زاده، نشر نی.
    \item اندرو هیوود، \textit{ایدئولوژی‌های سیاسی}، ترجمه محمد رفیعی مهرآبادی، نشر نی.
    \item جان اسکات، \textit{بنیان‌گذاران جامعه‌شناسی}، ترجمه جواد افشار، نشر دانشگاهی.
\end{itemize}

\vspace{15pt}

\section*{English References}

\begin{latin}
\begin{itemize}[label=--]
    \item Andrew Heywood, \textit{Political Ideologies: An Introduction}, Palgrave Macmillan, 2017.
    \item Friedrich Hayek, \textit{The Road to Serfdom}, University of Chicago Press, 1944.
    \item Karl Marx \& Friedrich Engels, \textit{The Communist Manifesto}, 1848.
    \item John Locke, \textit{Two Treatises of Government}, 1689.
    \item Adam Smith, \textit{The Wealth of Nations}, 1776.
    \item Robert Nozick, \textit{Anarchy, State, and Utopia}, Basic Books, 1974.
    \item John Rawls, \textit{A Theory of Justice}, Harvard University Press, 1971.
\end{itemize}
\end{latin}

\vspace{20pt}

\begin{center}
\rule{0.4\textwidth}{0.5pt}
\end{center}
